\naslov{Linearna regresija}

Linearna regresija je model dogajanja, kjer pojasnimo opažene podatke kot
$\sv{Y} = X \sv{\beta} + \sv{\varepsilon}$ za konstantno matriko opažanj $X$,
determinističen (a neznan) vektor $\sv{\beta}$ in slučajni šum
$\sv{\varepsilon}$.
Predpostavimo, da je $\sv{\varepsilon}$ \pojem{beli šum}, torej da je
$E(\sv{\varepsilon}) = \sv{0}$ in $\var(\sv{\varepsilon}) = \sigma^2 I$.
Druga lastnost je kombinacija nekoreliranosti komponent in
\pojem{homoskedastičnosti}, da so diagonalni elementi enaki.
Rešitev za $\sv{\beta}$ po najmanjših kvadratih je $\hat{\sv{\beta}} = (X^T
X)^{-1} X^T \sv{Y}$.

\begin{definicija}
  \pojem{Linearna cenilka} za količino $q$ z vrednostmi v $\R$ na podlagi
  opažanja $\sv{Y} \in \R^n$ je cenilka oblike $\hat{q} = \sv{\psi}^T \sv{Y} +
  \rho$, kjer sta $\sv{\psi} \in \R^n$ in $\rho \in \R$ konstanti.
\end{definicija}

\begin{izrek}[Gauss, Markov]
  Če je $\hat{\sv{\beta}}$ cenilka po najmanjših kvadratih kot zgoraj, je
  $\sv{c}^T \hat{\sv{\beta}}$ najboljša nepristranska linearna cenilka za
  $\sv{c}^T \sv{\beta}$.
\end{izrek}

\begin{proof}
  Model lahko parametriziramo že z $\sv{v} = X \sv{\beta}$.
  Obstaja tak $\sv{\varphi} \in V = \im X$, da je $q = \sv{\varphi}^T \sv{v}$.
  Definiramo $\hat{\sv{v}} = X \hat{\sv{\beta}} = H \sv{Y}$ za $H = X (X^T
  X)^{-1} X^T$.
  Ekvivalentno je pokazati, da je $\sv{\varphi}^T \hat{\sv{v}}$ NNLC za $q =
  \sv{\varphi}^T \sv{v}$.

  Linearne cenilke za $q$ so oblike $\hat{Q} = \sv{\psi}^T \sv{Y} + \rho$.
  Cenilka je nepristranska, če je
  \[
	E(\hat{Q}) = \sv{\psi}^T E(\sv{Y}) + \rho
	= \sv{\psi}^T X \sv{\beta} + \rho = \sv{\psi}^T \sv{v} + \rho
	= \sv{\varphi}^T \sv{v}.
  \]
  To mora veljati za vsak $\sv{v}$.
  Pri $\sv{v} = 0$ dobimo $\rho = 0$, iz preostalih pa sklepamo:
  \[
	\sv{\psi}^T \sv{v}
	= \sv{\psi}^T H \sv{v}
	= (H \sv{\psi})^T \sv{v}
	= \sv{\varphi}^T \sv{v}
  \]
  upoštevaje dejstvo, da je $H$ simetričen projektor.
  Sledi $\sv{\psi} = \sv{\varphi} + \sv{\eta}$, kjer je $\sv{\eta} \in V^\bot$,
  torej $\hat{Q} = (\sv{\varphi} + \sv{\eta})^T \sv{Y}$.

  Če izračunamo varianco, dobimo
  \[
	\var(\hat{Q}) = (\sv{\varphi} + \sv{\eta})^T \var(\sv{\varepsilon})
	(\sv{\varphi} + \sv{\eta})
	= \sigma^2 (\sv{\varphi}^T \sv{\varphi} + \sv{\eta}^T \sv{\eta}),
  \]
  torej je varianca minimalna pri $\sv{\eta} = 0$.
\end{proof}

\vprasanje{Povej in dokaži izrek Gaussa in Markova.}

Za vsoto kvadratov ostankov
\[
  \text{RSS} = \norm{\sv{Y} - X \hat{\sv{\beta}}}^2
  = \norm{(I_n - H) \sv{Y}}^2
\]
lahko s pomočjo trika s sledjo izračunamo
\[
  E(\text{RSS}) = \sigma^2 (n - \rang H) = \sigma^2 (n-p),
\]
torej dobimo nepristransko cenilko $\hat{\sigma}^2_+ = \frac{\text{RSS}}{n-p}$
za $\sigma^2$.

\vprasanje{Izpelji nepristransko cenilko za varianco šuma pri linearni
  regresiji.}

Če vzamemo Gaussov model šuma, torej $\sv{\varepsilon} \sim N(\sv{0}, \sigma^2
I_n)$, potem je $\hat{\sv{\beta}}$ tudi cenilka za $\sv{\beta}$ po metodi
največjega verjetja.
Če je $\hat{\sigma}$ cenilka po MNV za $\sigma$, pa je $\hat{\sigma}^2 =
\frac{\text{RSS}}{n}$.
Izkaže se
\[
  \frac{\text{RSS}}{\sigma^2} \sim \chi^2(n-p)
\]
in
\[
  \frac{\sv{c}^T \hat{\sv{\beta}} - \sv{c}^T \sv{\beta}}{\hat{\sigma}^2_+
	\sqrt{\sv{c}^T (X^T X)^{-1} \sv{c}}}
  \sim \operatorname{Student}(n-p).
\]


% LocalWords:  nekoreliranosti homoskedastičnosti NNLC
